\documentclass[conference]{IEEEtran}
\usepackage{graphicx}
\usepackage{float}
\usepackage{placeins}
\usepackage{amsmath}
\affiliation{\institution{Open Research Lab}\country{USA}}
\email{jho@example.org}

\graphicspath{{./}}

\begin{document}
\title{Denoising Diffusion Models in Practice}
\author{J. Ho}
\maketitle
\begin{abstract}
Diffusion models iteratively denoise random noise into structured samples.
\end{abstract}
\begin{IEEEkeywords}
diffusion models, generative modeling
\end{IEEEkeywords}
\IEEEoverridecommandlockouts
\section{Introduction}
Likelihood-based generation can be stable and high quality.

\FloatBarrier
\section{Training}
The simplified loss is
\begin{equation}
\mathcal{L}_{simple}=\mathbb{E}_{t,x,\epsilon}\left[\|\epsilon-\epsilon_\theta(x_t,t)\|_2^2\right].
\end{equation}

\FloatBarrier
\section{Conclusion}
Strong guidance and schedulers improve final fidelity.

\FloatBarrier
\bibliographystyle{IEEEtran}
\bibliography{references}
\end{document}
